\chapter{Language Design}

The language design is intentionally small at the surface. A scalar binding
names a value; an indexed binding builds a tensor point by point; a reduction
introduces a local index; a recurrence defines a binding through base and
recurrent clauses; and a derivative request asks for the sensitivity of one
named value with respect to another.

\begin{lstlisting}[style=einlang]
let C[i, j] = sum[k](A[i, k] * B[k, j]);
let x[0] = x0;
let x[t in 1..T] = step(x[t - 1], u[t]);
let dx0 = @x[T - 1] / @x0;
\end{lstlisting}

The matrix multiplication line exposes output axes \texttt{i} and \texttt{j}
and reduction axis \texttt{k}. The recurrence exposes the self-read offset
\texttt{t - 1}. The derivative request names the target and parameter at the
point where the sensitivity is needed.

The example is intentionally small, but it shows the thesis boundary. The
compiler does not receive an opaque tensor call followed by a host loop followed
by an AD transform; it receives three source constructs in one binding
environment. That means the name resolver, range analysis, recurrence analysis,
and autodiff passes can agree on the same DefIds. The programmer still writes
mathematical notation, while the implementation keeps enough structure to make
compile-time checks and lowering choices before execution.

\section{Indexed Definitions}

An indexed binding defines a tensor over an index environment. Explicit ranges
such as \texttt{i in 0..M} attach domains directly to variables. When ranges are
omitted, shaped tensor reads can provide candidate domains: a read of
\texttt{A[i,k]} constrains \texttt{i} to the first axis of \texttt{A} and
\texttt{k} to the second. Reduction indices follow the same rule inside the
reduction body.

\begin{lstlisting}[style=einlang]
let Y[n in 0..N, o in 0..OC, i in 0..OH, j in 0..OW] =
    sum[c in 0..IC, r in 0..KH, s in 0..KW](
        X[n, c, i + r, j + s] * W[o, c, r, s]);
\end{lstlisting}

This direct convolution keeps surviving axes, reduction axes, and affine input
offsets in one source expression. An \texttt{einsum}-plus-patches encoding can
run the same computation, but it splits the sliding-window construction from the
contraction. \einlang{} keeps both facts in one binding.

The analysis value is the combination, not only the notation. The output axes
\texttt{n}, \texttt{o}, \texttt{i}, and \texttt{j} define the produced tensor,
while the reduction axes \texttt{c}, \texttt{r}, and \texttt{s} define the
contracted kernel domain. The affine offsets \texttt{i + r} and \texttt{j + s}
also remain visible to bounds checking and lowering. A backend may still choose
to run this as patches, scalar loops, or a vectorized contraction, but that is a
later execution decision rather than the source-level representation.

\begin{center}
\centering
\begin{minipage}[t]{0.48\linewidth}
\textbf{Einlang: named axes in IR}
\begin{lstlisting}[style=einlang,basicstyle=\ttfamily\footnotesize]
let C[i, j] =
    sum[k](A[i, k] * B[k, j]);
\end{lstlisting}
\end{minipage}
\hfill
\begin{minipage}[t]{0.48\linewidth}
\textbf{NumPy: string mini-language}
\begin{lstlisting}[style=pythoncode,basicstyle=\ttfamily\footnotesize]
C = np.einsum("ik,kj->ij", A, B)
\end{lstlisting}
\end{minipage}

\vspace{0.45\baselineskip}
\begin{minipage}[t]{0.48\linewidth}
\textbf{TVM: tensor expression}
\begin{lstlisting}[style=pythoncode,basicstyle=\ttfamily\footnotesize]
k = te.reduce_axis((0, K), "k")
C = te.compute((M, N),
    lambda i, j: te.sum(
        A[i, k] * B[k, j], axis=k))
\end{lstlisting}
\end{minipage}
\hfill
\begin{minipage}[t]{0.48\linewidth}
\textbf{Halide: update definition}
\begin{lstlisting}[style=pythoncode,basicstyle=\ttfamily\footnotesize]
Var i, j;
RDom k(0, K);
Func C;
C(i, j) += A(i, k) * B(k, j);
\end{lstlisting}
\end{minipage}

\captionsetup{type=figure,hypcap=false}
\caption{Thesis-adapted comparison of indexed reduction interfaces. The reused
paper example is copied into the thesis figure set and retuned around compiler
facts rather than presentation brevity.}
\label{fig:thesis-indexed-reduction-interfaces}
\end{center}

Figure~\ref{fig:thesis-indexed-reduction-interfaces} should be read as a
compiler-boundary comparison. NumPy places the contraction structure in a
string, which is compact but separated from ordinary name resolution. TVM and
Halide expose structured tensor expressions, but they are embedded in host-side
scheduling systems. The \einlang{} form keeps the same source binding, the same
index variables, and the same diagnostics context available to later passes.
That is why the thesis treats indexed notation as an implementation contract
rather than as surface sugar.

\subsection{Index Domains and Where Clauses}

Index domains are intentionally attached to index positions rather than hidden
inside a separate loop statement. This matters for diagnostics and lowering. A
range such as \texttt{i in 0..M} is in scope for the indexed body, and a
reduction range such as \texttt{k in 0..K} is in scope only for the reduction
body. Predicate-style \texttt{where} clauses refine a domain, while iteration
ranges in \texttt{where} clauses are rejected by validation because they obscure
which source construct owns the index.

\begin{lstlisting}[style=einlang]
let even_square[i in 0..N] = i * i where i % 2 == 0;
let total = sum[k in 0..N](even_square[k]);
\end{lstlisting}

The implementation uses this discipline in range analysis: every index variable
has a DefId, and all inferred domains are keyed by that DefId. The same spelling
can be reused in separate indexed bindings without aliasing compiler state.

The two-line example demonstrates both parts of the rule. The first binding
creates a tensor only over even points in \texttt{0..N}; the second reduction
then consumes that tensor through an ordinary loop variable. Because the
\texttt{where} predicate is not itself a new iteration domain, the compiler can
classify it as a guard. That distinction matters later: guards become lowered
conditions, while ranges become loop bounds. Collapsing the two would make it
harder to explain diagnostics and harder for the backend to choose a correct
execution strategy.

\section{Declarative Recurrence}

A recurrent binding may have multiple clauses. Base clauses define boundary
points; recurrent clauses define the remaining domain and may read other points
of the same binding.

\begin{lstlisting}[style=einlang]
let H[0, j in 0..N] = top[j];
let H[i in 0..M, 0] = left[i];
let H[i in 1..M, j in 1..N] =
    mix(H[i - 1, j], H[i, j - 1], input[i, j]);
\end{lstlisting}

The source contract is dependency, not schedule. The compiler must evaluate each
point after the points it reads, but the source program does not have to pick a
loop nest, carry tuple, row buffer, or full materialization policy up front.

This example is a two-dimensional dependency relation. The interior clause
reads one point above and one point to the left, so a legal traversal must
respect both predecessor directions. The source does not say whether rows,
columns, tiles, or a full materialized grid should be used; it states the
relation that any such schedule must satisfy. The implementation can therefore
derive recurrence dimensions and storage metadata from reads of the same DefId,
instead of attempting to reverse-engineer a relation from already-written host
control flow.

\begin{center}
\centering
\begin{minipage}[t]{0.48\linewidth}
\textbf{Einlang: dependency relation}
\begin{lstlisting}[style=einlang,basicstyle=\ttfamily\footnotesize]
let y[0] = init;
let y[t in 1..T] =
    step(y[t - 1], x[t]);
let out = y[T - 1];
\end{lstlisting}
\end{minipage}
\hfill
\begin{minipage}[t]{0.48\linewidth}
\textbf{JAX: carry contract}
\begin{lstlisting}[style=pythoncode,basicstyle=\ttfamily\footnotesize]
def body(carry, x_t):
    y = step(carry, x_t)
    return y, y
out, ys = jax.lax.scan(
    body, init, x[1:])
\end{lstlisting}
\end{minipage}

\vspace{0.45\baselineskip}
\begin{minipage}[t]{0.48\linewidth}
\textbf{PyTorch: host storage}
\begin{lstlisting}[style=pythoncode,basicstyle=\ttfamily\footnotesize]
ys = [init]
for x_t in x[1:]:
    ys.append(step(ys[-1], x_t))
out = ys[-1]
\end{lstlisting}
\end{minipage}
\hfill
\begin{minipage}[t]{0.48\linewidth}
\textbf{Compiler-visible fact}
\begin{lstlisting}[style=pythoncode,basicstyle=\ttfamily\footnotesize]
self_reads = {t: t - 1}
order = topological_order(self_reads)
storage = bounded_window(self_reads)
\end{lstlisting}
\end{minipage}

\captionsetup{type=figure,hypcap=false}
\caption{Thesis-adapted comparison of recurrence interfaces. The thesis copy
emphasizes which scheduling and storage decisions are source obligations in
neighboring systems and which are compiler obligations in \einlang{}.}
\label{fig:thesis-recurrence-interfaces}
\end{center}

Figure~\ref{fig:thesis-recurrence-interfaces} highlights the delayed
commitment. JAX asks the user to choose a carry boundary and stacked output
shape before the compiler sees later observations. A PyTorch loop makes storage
an ordinary host-language list decision. In \einlang{}, the self-read offset is
the source fact; traversal order and finite-window storage are compiler
obligations. This difference is small in a one-step recurrence, but it becomes
important when later code observes only a suffix, a final value, or an
intermediate derivative.

\subsection{Boundary Clauses}

Base and recurrent clauses share one binding name. The grouping pass therefore
sees the following as one recurrent definition rather than as unrelated
assignments:

\begin{lstlisting}[style=einlang]
let fib[0] = 0;
let fib[1] = 1;
let fib[n in 2..N] = fib[n - 1] + fib[n - 2];
\end{lstlisting}

This source form is small, but it carries three compiler facts: the covered
boundary points, the recurrent domain, and the static self-read offsets. Those
facts are consumed later by lowering, recurrence-order analysis, and storage
planning.

The base clauses are not special cases outside the binding. They are part of
the same grouped definition as the recurrent clause, which lets the compiler
check that the values read by \texttt{fib[n - 1]} and \texttt{fib[n - 2]} are
available when the recurrent body runs. The same grouping also prevents the
runtime from treating the three statements as independent outputs. The
important analysis point is that boundary coverage and recurrence dependency
are represented together, so storage inference can distinguish required
history from values that only initialize the sequence.

\section{Local Derivative Requests}

Derivative requests are expressions over named source bindings:

\begin{lstlisting}[style=einlang]
let loss = sum[i]((pred[i] - target[i]) * (pred[i] - target[i]));
let dloss_dW = @loss / @W;
\end{lstlisting}

The target and parameter names determine the derivative shape. Because the
request is an expression, it can appear near the computation it differentiates,
including near intermediate or recurrent values. This is narrower than a full
host-language AD system, but it gives the source language a simple and local
differentiation interface.

The quotient form is restricted to named source bindings because this makes the
compiler's dependency query precise. A derivative of an anonymous inline
subexpression would force the compiler to invent a hidden binding and decide its
scope. \einlang{} instead asks the programmer to name the value, which also
improves diagnostics and reuse.

The two-line derivative example is deliberately source-local. The user names
the loss once, then asks for the sensitivity with respect to \texttt{W} where
the result is needed. The compiler can use the DefId of \texttt{loss} as the
target and the DefId of \texttt{W} as the parameter, so the request is a
structured dependency query rather than a string or callback. This is also why
anonymous derivatives are not the core interface: naming the target gives the
implementation a stable node to analyze, cache, lower, and report in
diagnostics.

\begin{center}
\centering
\begin{minipage}[t]{0.48\linewidth}
\textbf{Einlang: shaped source expression}
\begin{lstlisting}[style=einlang,basicstyle=\ttfamily\footnotesize]
let local[t, j] =
    @h[t, j] / @theta[j];
\end{lstlisting}
\end{minipage}
\hfill
\begin{minipage}[t]{0.48\linewidth}
\textbf{JAX: package selected value}
\begin{lstlisting}[style=pythoncode,basicstyle=\ttfamily\footnotesize]
local = jax.grad(
    lambda th: run(th)[t, j]
)(theta)[j]
\end{lstlisting}
\end{minipage}

\vspace{0.45\baselineskip}
\begin{minipage}[t]{0.48\linewidth}
\textbf{Julia/Zygote: closure boundary}
\begin{lstlisting}[style=pythoncode,basicstyle=\ttfamily\footnotesize]
using Zygote
local = gradient(
    th -> run(th)[t, j], theta
)[1][j]
\end{lstlisting}
\end{minipage}
\hfill
\begin{minipage}[t]{0.48\linewidth}
\textbf{PyTorch: graph-state query}
\begin{lstlisting}[style=pythoncode,basicstyle=\ttfamily\footnotesize]
h = run(theta)
theta.grad = None
h[t, j].backward()
local = theta.grad[j]
\end{lstlisting}
\end{minipage}

\captionsetup{type=figure,hypcap=false}
\caption{Thesis-adapted comparison of local derivative interfaces. The thesis
copy keeps the paper's motivating query but retunes the labels around the
implementation boundary that the autodiff passes consume.}
\label{fig:thesis-local-derivative-interfaces}
\end{center}

Figure~\ref{fig:thesis-local-derivative-interfaces} compares where the
derivative boundary appears. JAX and Zygote package the selected value behind a
function-like object; PyTorch exposes the query through graph state and a
gradient buffer. \einlang{} instead makes the sensitivity an expression whose
shape follows from the target and parameter. The implementation burden is
higher because the compiler must lower the shaped request, but the source
benefit is that intermediate and recurrent sensitivities stay in the same
indexed language as the values they describe.

\section{Types and Values}

The executed implementation supports primitive scalars, rectangular arrays,
jagged arrays for non-Einstein data, tuples, function values, and differential
types used during the autodiff phase. Rectangular arrays are the tensor type
that participates in indexed notation. A type can specify concrete dimensions,
unknown dimensions, or dynamic rank:

\begin{lstlisting}[style=einlang]
let x: f32 = 1.0;
let v: [f32] = [1.0, 2.0, 3.0];
let m: [f32; 2, 3] = [[1.0, 2.0, 3.0], [4.0, 5.0, 6.0]];
let dyn: [f32; *] = load_npy("weights.npy");
\end{lstlisting}

Shape is not only a type annotation. The compiler separately computes
expression-level shape facts from literals, indexed definitions, reductions,
and calls. This separation lets the type system track precision and rank while
shape analysis tracks the concrete dimensions needed for indexing and lowering.

The example distinguishes four levels of precision. A scalar has a primitive
type only; a vector fixes rank but may leave length dynamic; a rectangular
literal can expose concrete dimensions; and a loaded array may have dynamic
rank. The compiler keeps these cases separate so that a function can be type
compatible even when a later range or lowering pass still needs a concrete
shape fact. In other words, type checking decides whether an expression is a
tensor-like value of the right element kind, while shape analysis decides which
indices and loops are legal for that value.

\section{Functions, Blocks, and Modules}

Functions are expressions with parameters, optional type annotations, and an
implicit return value equal to the final block expression. Functions can be
called before their textual definition after name resolution has hoisted item
bindings. Modules are file-backed units resolved through \texttt{use} and
\texttt{mod} declarations. Public functions and public re-exports make the
standard library usable as ordinary \einlang{} code rather than as a bag of
compiler intrinsics.

\begin{lstlisting}[style=einlang]
use std::math::basic::sqrt;

fn norm2(x, y) {
    sqrt(x * x + y * y)
}

let r = norm2(3.0, 4.0);
\end{lstlisting}

This snippet shows the ordinary-language machinery that the tensor features
depend on. The \texttt{use} statement must resolve through the module system
before the function call can be typed, and the function body must be hoisted so
the call site can refer to a stable DefId. The final expression of the block is
the return value, which keeps functions lightweight enough to use in indexed
bodies. Without this module and function layer, the tensor core would collapse
into a set of builtins rather than a language with reusable source libraries.

The example also shows why ordinary expression features are part of the
implementation story rather than background convenience. A standard-library
function such as \texttt{sqrt} has to be imported, assigned a DefId, checked as
a callable value, and executed through the same backend path as user functions.
The final-expression return rule lets helper functions remain small enough to
appear naturally inside tensor bodies, while the module system lets those
helpers live outside the compiler. This is what keeps tensor notation from
becoming a closed DSL with only hard-coded operations.

\section{General Language Surface}

The implementation supports more than the tensor core. The grammar includes
function definitions, public functions, \texttt{@fn} differential rules, typed
parameters and returns, immutable \texttt{let} and \texttt{const} declarations,
tuple destructuring, blocks as expressions, \texttt{if}, \texttt{match},
patterns, arrays, tuples, member access, pipeline expressions, casts,
interpolated strings, module paths, \texttt{use} statements, inline modules,
struct and enum syntax, and where clauses. Some parsed surface forms are
implemented as current language features, while some remain parsed but not yet
fully supported by the backend. The thesis focuses on the non-trivial executed
path: tensor definitions, functions, modules, typing, recurrences, autodiff,
standard-library calls, and examples.

\section{Composition}

The constructs are designed to be orthogonal. An indexed definition can feed a
recurrence; a recurrence can feed an indexed reduction; a derivative request can
target a final scalar or an intermediate tensor. The language benefit is stable
notation. The compiler benefit is that the same IR can carry axes, domains,
dependency offsets, and derivative targets until the pass that consumes each
fact has run.

\begin{lstlisting}[style=einlang]
let score[t] = sum[d](x[t, d] * w[d]);
let prefix[0] = score[0];
let prefix[t in 1..T] = prefix[t - 1] + score[t];
let dw = @prefix[T - 1] / @w;
\end{lstlisting}

The same source region contains a contraction, a recurrent binding, and a local
derivative request. In a host framework this would typically become an array
operation, a loop or scan, and a separate AD call. In \einlang{}, the compiler
sees the named axes, the recurrent dependency, and the derivative target before
any of them are lowered to backend operations.

Composition is the strongest reason to keep the constructs in one language.
The \texttt{score} binding exposes an indexed contraction, \texttt{prefix}
turns that value into a recurrent sequence, and \texttt{dw} asks for a
derivative of the final recurrent value. Each line consumes facts produced by
earlier lines without changing notation families. The compiler can therefore
thread range facts, recurrence offsets, and autodiff dependencies through one
IR instead of coordinating several external APIs. This is the source-level
version of the implementation thesis: smart passes become possible because the
program has not hidden its structure too early.
